\documentclass[10pt,a4paper]{article}

% ==============================================================================
% EXPERIMENT 01: PRINTABLE CODE, PROTOCOLS & OUTPUT (EXACT SINGLE 1-PAGE DOCUMENT)
% Formatted with 1.0-inch lateral margins, large readable fonts, and tight top/bottom margins.
% ==============================================================================
\usepackage[utf8]{inputenc}
\usepackage[T1]{fontenc}
\usepackage{lmodern}
\usepackage[left=1.0in, right=1.0in, top=0.35in, bottom=0.35in]{geometry}
\usepackage{amsmath,amssymb}
\usepackage{listings}
\usepackage{tikz}
\usetikzlibrary{shapes.geometric,arrows.meta,positioning,calc}
\usepackage{booktabs}
\usepackage{tabularx}
\usepackage{titlesec}
\usepackage{caption}
\usepackage{microtype}
\usepackage{float}

% Disable all running headers, footers, and page numbers for clean print & paste
\pagestyle{empty}
\setlength{\parindent}{0pt}

% Section Titles Formatting (Enlarged)
\titleformat{\section}
  {\normalfont\normalsize\bfseries\raggedright}{\thesection}{0.5em}{}[\titlerule]

\titlespacing*{\section}{0pt}{5pt}{2pt}

% Code Listing Styling (Strict Monochrome with Enlarged Font for High-Contrast Print)
\lstdefinestyle{bwCode}{
    language=C++,
    basicstyle=\ttfamily\scriptsize,
    keywordstyle=\bfseries,
    commentstyle=\itshape,
    stringstyle=\ttfamily,
    numbers=left,
    numberstyle=\tiny,
    numbersep=5pt,
    stepnumber=1,
    breaklines=true,
    breakatwhitespace=true,
    frame=single,
    rulecolor=\color{black},
    tabsize=2,
    showstringspaces=false,
    columns=fullflexible,
    keepspaces=true,
    lineskip=-0.2pt,
    aboveskip=2pt,
    belowskip=2pt
}

\lstdefinestyle{bwTerminal}{
    basicstyle=\ttfamily\scriptsize,
    numbers=none,
    breaklines=true,
    breakatwhitespace=true,
    frame=single,
    rulecolor=\color{black},
    columns=fullflexible,
    keepspaces=true,
    lineskip=0.0pt,
    aboveskip=2pt,
    belowskip=2pt
}

\captionsetup{font=footnotesize,labelfont=bf,skip=2pt}

\begin{document}

% ==============================================================================
% DOCUMENT HEADER
% ==============================================================================
\begin{center}
    {\large\bfseries Experiment 01: Overview of IoT Architecture \& Protocols}\\[0.06cm]
    {\normalsize\textbf{5-Tier Architecture, Protocol Comparison \& MQTT Telemetry Payload Code}}
\end{center}
\vspace{-0.24cm}
\rule{\textwidth}{0.8pt}
\vspace{0.06cm}

% ==============================================================================
% 1. PROTOCOL COMPARATIVE MATRIX
% ==============================================================================
\section*{1. Constrained IoT Application Protocols Matrix}

\begin{table}[H]
\centering
\scriptsize
\setlength{\tabcolsep}{4.5pt}
\renewcommand{\arraystretch}{1.20}
\begin{tabularx}{\textwidth}{|p{1.8cm}|p{1.7cm}|p{3.2cm}|p{2.1cm}|X|}
\hline
\textbf{Protocol} & \textbf{Transport} & \textbf{Messaging Model} & \textbf{Header Size} & \textbf{Primary Application Domain} \\
\hline
\textbf{MQTT} & TCP (1883) & Publish / Subscribe & 2 Bytes (Fixed) & Real-time sensor telemetry, SCADA, battery-backed nodes. \\
\textbf{CoAP} & UDP (5683) & REST Request / Response & 4 Bytes (Fixed) & Constrained 8-bit MCUs, low-power lossy networks (LLNs). \\
\textbf{HTTP/REST} & TCP (80/443) & Client / Server Request & $>500$ B (Text) & Cloud-to-cloud APIs, enterprise REST integrations. \\
\textbf{WebSocket} & TCP (80) & Full-Duplex Stream & 2--10 Bytes & Live browser telemetry monitoring dashboards. \\
\hline
\end{tabularx}
\caption{Technical Comparison of IoT Application Layer Protocols.}
\end{table}

\vspace{-0.1cm}

% ==============================================================================
% 2. MQTT TELEMETRY SERIALIZER SOURCE CODE
% ==============================================================================
\section*{2. Edge Microcontroller Telemetry Payload Firmware}

\begin{lstlisting}[style=bwCode, title={\textbf{Listing 1.1: Microcontroller Firmware for MQTT Telemetry Serialization.}}]
/*
 * Experiment 01: IoT Edge Telemetry Serialization & Protocol Ingestion
 * Formats multi-sensor readings into lightweight JSON payloads for MQTT Broker publishing.
 */

#include <ArduinoJson.h> // Standard lightweight embedded JSON library

void setup() {
  Serial.begin(9600); // Initialize UART Serial Interface
  Serial.println("==================================================");
  Serial.println("   IoT Lab: MQTT Node Telemetry Payload Stream   ");
  Serial.println("==================================================");
}

void loop() {
  // Step 1: Acquire mock sensory inputs from Perception Layer
  float currentTemp = 28.4;
  float currentHumidity = 62.5;
  int deviceUptime = millis() / 1000;

  // Step 2: Construct serialized JSON payload packet
  StaticJsonDocument<128> doc;
  doc["node_id"] = "ESP32_NODE_01";
  doc["uptime"]  = deviceUptime;
  doc["temp"]    = currentTemp;
  doc["hum"]     = currentHumidity;
  doc["status"]  = "ONLINE";

  // Step 3: Stream MQTT packet to Serial Gateway
  Serial.print("TOPIC: sensors/room1/telemetry | PAYLOAD: ");
  serializeJson(doc, Serial);
  Serial.println();

  delay(3000); // 3-second transmission cycle
}
\end{lstlisting}

\vspace{-0.1cm}

% ==============================================================================
% 3. TELEMETRY OUTPUT & BROKER LOG
% ==============================================================================
\section*{3. Simulated MQTT Broker Telemetry Stream Output}

\begin{lstlisting}[style=bwTerminal]
==================================================
   IoT Lab: MQTT Node Telemetry Payload Stream   
==================================================
TOPIC: sensors/room1/telemetry | PAYLOAD: {"node_id":"ESP32_NODE_01","uptime":3,"temp":28.4,"hum":62.5,"status":"ONLINE"}
TOPIC: sensors/room1/telemetry | PAYLOAD: {"node_id":"ESP32_NODE_01","uptime":6,"temp":28.6,"hum":62.1,"status":"ONLINE"}
TOPIC: sensors/room1/telemetry | PAYLOAD: {"node_id":"ESP32_NODE_01","uptime":9,"temp":28.9,"hum":61.8,"status":"ONLINE"}
TOPIC: sensors/room1/telemetry | PAYLOAD: {"node_id":"ESP32_NODE_01","uptime":12,"temp":29.1,"hum":61.5,"status":"ONLINE"}
\end{lstlisting}

\end{document}
